\documentclass[aspectratio=169,11pt]{beamer}

\usepackage[T1]{fontenc}
\usepackage[utf8]{inputenc}
\usepackage[brazilian]{babel}
\usepackage{lmodern}
\usepackage{amsmath,amssymb}
\usepackage{booktabs,tabularx,array}
\usepackage{xcolor}
\usepackage{tikz}
\usepackage{pgfplots}
\usepackage{listings}
\usepackage{ragged2e}
\usepackage{hyperref}
\usetikzlibrary{positioning,arrows.meta,shapes.geometric,calc,decorations.pathreplacing}
\pgfplotsset{compat=1.18}

\definecolor{MondrianRed}{HTML}{C83B32}
\definecolor{MondrianBlue}{HTML}{315B91}
\definecolor{MondrianYellow}{HTML}{E8BD45}
\definecolor{MondrianBlack}{HTML}{171717}
\definecolor{SoftGray}{HTML}{F3F4F2}
\definecolor{MidGray}{HTML}{686B70}

\setbeamercolor{normal text}{fg=MondrianBlack,bg=white}
\setbeamercolor{structure}{fg=MondrianBlue}
\setbeamercolor{frametitle}{fg=MondrianBlack,bg=white}
\setbeamercolor{title}{fg=MondrianBlack}
\setbeamercolor{block title}{fg=MondrianBlack,bg=SoftGray}
\setbeamercolor{block body}{fg=MondrianBlack,bg=SoftGray}
\setbeamercolor{alerted text}{fg=MondrianRed}
\setbeamerfont{frametitle}{series=\bfseries,size=\Large}
\setbeamertemplate{navigation symbols}{}
\setbeamertemplate{itemize item}{\color{MondrianBlue}\small$\blacksquare$}
\setbeamertemplate{itemize subitem}{\color{MondrianRed}\scriptsize$\blacksquare$}
\setbeamertemplate{footline}{
  \leavevmode
  \begin{beamercolorbox}[wd=\paperwidth,ht=2.6ex,dp=1.1ex,sep=0pt]{normal text}
    \hspace{1em}\color{MidGray}\scriptsize Quick Sort \hspace{.5em}\textbar\hspace{.5em} Estrutura de Dados \hspace{.5em}\textbar\hspace{.5em} UFRR
    \hfill\color{MidGray}\scriptsize\insertframenumber\,/\,\inserttotalframenumber\hspace{1em}
  \end{beamercolorbox}
}
\setbeamertemplate{frametitle}{
  \vspace{.25em}
  \insertframetitle\par
  \vspace{-.45em}
  {\color{MondrianBlack}\rule{\textwidth}{.45pt}}
}
\hypersetup{colorlinks=true,linkcolor=MondrianBlue,urlcolor=MondrianBlue,citecolor=MondrianBlue}

\lstdefinestyle{code}{
  language=C,
  basicstyle=\ttfamily\scriptsize,
  keywordstyle=\bfseries\color{MondrianBlue},
  commentstyle=\itshape\color{MidGray},
  stringstyle=\color{MondrianRed},
  numbers=left,
  numberstyle=\tiny\color{MidGray},
  numbersep=5pt,
  showstringspaces=false,
  breaklines=true,
  frame=single,
  rulecolor=\color{MondrianBlack},
  xleftmargin=1em,
  framexleftmargin=.7em,
  tabsize=4,
  columns=fullflexible
}

\newcommand{\citep}[1]{{\color{MidGray}\scriptsize\cite{#1}}}
\newcommand{\key}[1]{\textbf{\color{MondrianBlue}#1}}
\newcommand{\softbox}[2]{%
  \begin{beamercolorbox}[wd=\linewidth,sep=.8em,rounded=false]{block body}
    {\bfseries #1}\par\vspace{.25em}#2
  \end{beamercolorbox}
}
\newcommand{\placeholder}[2]{%
  \begin{tikzpicture}
    \node[draw=MondrianBlack,line width=.7pt,minimum width=#1,minimum height=#2,
      fill=SoftGray,align=center,text=MidGray] {\scriptsize Espaço para inserir figura\\[-1pt]
      \tiny (captura de tela / foto / imagem)};
  \end{tikzpicture}
}

\title[Quick Sort]{Quick Sort}
\subtitle{Ideia, análise, variantes e evidências experimentais}
\author{Nome do(a) aluno(a)}
\institute{Universidade Federal de Roraima\\Departamento de Ciência da Computação\\Estrutura de Dados\\Prof. Filipe Dwan Pereira}
\date{Data: \rule{2.1cm}{.4pt}}

\begin{document}

\begin{frame}[plain]
  \begin{tikzpicture}[remember picture,overlay]
    \fill[MondrianBlack] (current page.north west) rectangle ([xshift=.18cm]current page.south west);
    \fill[MondrianBlue] ([xshift=.42cm,yshift=-.5cm]current page.north west) rectangle ([xshift=1.25cm,yshift=-2.4cm]current page.north west);
    \fill[MondrianRed] ([xshift=1.43cm,yshift=-.5cm]current page.north west) rectangle ([xshift=2.15cm,yshift=-1.42cm]current page.north west);
    \fill[MondrianYellow] ([xshift=.42cm,yshift=-2.6cm]current page.north west) rectangle ([xshift=1.25cm,yshift=-3.55cm]current page.north west);
    \draw[MondrianBlack,line width=.8pt] ([xshift=.32cm,yshift=-.4cm]current page.north west) rectangle ([xshift=2.25cm,yshift=-3.7cm]current page.north west);
  \end{tikzpicture}
  \vspace{1.35cm}
  \hspace{2.15cm}
  \begin{minipage}{.75\textwidth}
    {\color{MidGray}\small UNIVERSIDADE FEDERAL DE RORAIMA \textbar\ DCC}\par
    \vspace{.55cm}
    {\Huge\bfseries Quick Sort}\par
    \vspace{.25cm}
    {\large Da partição à engenharia de ordenação}\par
    \vspace{.8cm}
    {\normalsize Estrutura de Dados \textbar\ Prof. Filipe Dwan Pereira}\par
    \vspace{.55cm}
    {\small Nome: \rule{5cm}{.4pt}\hspace{1em} Data: \rule{2cm}{.4pt}}
    \vfill
    {\scriptsize Uma leitura algorítmica, histórica e experimental.}
  \end{minipage}
\end{frame}

\begin{frame}{Roteiro e pergunta central}
  \begin{columns}[T,onlytextwidth]
    \column{.57\textwidth}
      \begin{enumerate}
        \item O que caracteriza o Quick Sort e como surgiu?
        \item Como a partição produz subproblemas menores?
        \item O que acontece, passo a passo, com seis valores?
        \item Por que o desempenho médio é bom, mas o pior caso é quadrático?
        \item O que a implementação e os dados da \texttt{lesson11} permitem concluir?
      \end{enumerate}
    \column{.38\textwidth}
      \softbox{Ideia em uma frase}{Escolher um pivô, particionar a sequência em torno dele e ordenar recursivamente os lados.}
      \vspace{.15em}
      \small
      \textbf{Objetivo:} explicar a técnica com precisão e distinguir garantias teóricas de medições de uma execução.
  \end{columns}
\end{frame}

\begin{frame}{O problema de ordenação}
  \begin{block}{Definição}
    Dada uma sequência de chaves \(A=(a_0,\ldots,a_{n-1})\), produzir uma permutação \(A'\) tal que
    \[
      a'_0\leq a'_1\leq\cdots\leq a'_{n-1}.
    \]
    O conteúdo deve ser preservado; muda-se a ordem, não o multiconjunto de valores.
  \end{block}
  \vspace{.4em}
  \begin{columns}[T,onlytextwidth]
    \column{.47\textwidth}
      \textbf{Propriedade de ordenação}\\
      Cada elemento aparece na posição compatível com a relação de ordem.
    \column{.47\textwidth}
      \textbf{Propriedade de permutação}\\
      Nenhuma chave é perdida, criada ou duplicada durante o processo.
  \end{columns}
  \vfill
  \citep{cormen2022,knuth1998}
\end{frame}

\begin{frame}{Contexto histórico: uma ideia que virou família}
  \begin{center}
    \begin{tikzpicture}[x=1cm,y=1cm]
      \draw[MondrianBlack,line width=1pt,-{Stealth[length=2mm]}] (0,0) -- (12.2,0);
      \foreach \x/\yr/\txt/\col in {
        1/1960/{Hoare concebe o método}/MondrianBlue,
        3.8/1962/{Artigo ``Quicksort''}/MondrianRed,
        6.7/1976/{Dijkstra: bandeira holandesa}/MondrianYellow,
        9.4/1993/{Bentley--McIlroy: engenharia}/MondrianBlue,
        11.4/1997/{Musser: introsort}/MondrianRed%
      } {
        \draw[MondrianBlack] (\x,0) -- (\x,.38);
        \fill[\col] (\x-.075,.32) rectangle (\x+.075,.47);
        \node[align=center,text width=2.2cm,font=\scriptsize] at (\x,1.05) {\txt};
        \node[font=\bfseries\scriptsize] at (\x,-.35) {\yr};
      }
    \end{tikzpicture}
  \end{center}
  \vspace{.2em}
  \begin{itemize}
    \item C. A. R. Hoare concebeu o algoritmo por volta de 1960; o artigo seminal saiu em 1962.
    \item O desenvolvimento posterior não é um único algoritmo: há diferentes partições, seleções de pivô e estratégias para casos especiais.
  \end{itemize}
  \citep{hoare1962,sedgwick2011,dijkstra1976,bentley1993,musser1997}
\end{frame}

\begin{frame}{O que significa particionar?}
  \begin{columns}[T,onlytextwidth]
    \column{.59\textwidth}
      Para um subvetor e pivô \(p\), uma partição correta estabelece:
      \[
      \underbrace{A[lo..q-1]}_{\leq p}\quad
      \underbrace{A[q]}_{p}\quad
      \underbrace{A[q+1..hi-1]}_{\geq p}
      \]
      \begin{itemize}
        \item O pivô termina em uma posição definitiva \(q\).
        \item Os dois lados ainda podem estar desordenados.
        \item Ordenam-se os lados; não é necessário juntar (merge) depois.
      \end{itemize}
    \column{.37\textwidth}
      \begin{tikzpicture}[x=.68cm,y=.68cm]
        \foreach \x/\v in {0/2,1/1,2/3,3/9,4/6,5/5} {
          \draw[MondrianBlack] (\x,0) rectangle ++(.86,.7);
          \node at (\x+.43,.35) {\small \v};
        }
        \draw[MondrianBlue,very thick] (0,-.18) -- (2.86,-.18);
        \draw[MondrianBlack,very thick] (2.86,-.18) -- (3.86,-.18);
        \draw[MondrianRed,very thick] (3.86,-.18) -- (5.86,-.18);
        \node[font=\scriptsize,align=center] at (1.4,-.68) {menores};
        \node[font=\scriptsize,align=center] at (3.35,-.68) {pivô};
        \node[font=\scriptsize,align=center] at (4.85,-.68) {maiores};
      \end{tikzpicture}
      \vspace{.5em}
      \small As relações exatas (\(<\), \(\leq\), etc.) dependem do esquema de partição e do tratamento de chaves iguais.
  \end{columns}
  \vfill
  \citep{cormen2022,sedgwick2011}
\end{frame}

\begin{frame}{Dividir e conquistar: recorrência e custo}
  \begin{columns}[T,onlytextwidth]
    \column{.53\textwidth}
      Se a partição deixa \(k\) elementos à esquerda:
      \[
        T(n)=T(k)+T(n-k-1)+\Theta(n).
      \]
      \begin{itemize}
        \item Partições equilibradas: \(T(n)=\Theta(n\log n)\).
        \item Partições extremas: \(T(n)=\Theta(n^2)\).
        \item Para chaves distintas e pivô aleatório, o número esperado de comparações é \(2(n+1)H_n-4n\), aproximadamente \(2n\ln n\).
      \end{itemize}
    \column{.43\textwidth}
      \softbox{Leitura da recursão}{Cada chamada paga pelo percurso da partição e depois resolve dois intervalos. O pivô é retirado dos subproblemas.}
      \vspace{.45em}
      \small
      \textbf{Espaço auxiliar:} \(O(1)\) para particionar in-place; além disso, a pilha recursiva é \(O(\log n)\) em divisões equilibradas e \(O(n)\) no pior caso.
  \end{columns}
  \vfill
  \citep{cormen2022,sedgwick2011}
\end{frame}

\begin{frame}{Um detalhe importante: Hoare e Lomuto não são a mesma partição}
  \footnotesize
  \begin{tabularx}{\textwidth}{>{\bfseries}p{2.5cm} X X}
    \toprule
    & Esquema de Hoare & Esquema Lomuto-style \\
    \midrule
    Movimento & Dois índices avançam de lados opostos. & Um índice de varredura e uma fronteira para os menores. \\
    Pivô & Nem sempre retorna já na posição final. & O pivô final é trocado para a fronteira. \\
    Recursão & Respeita a fronteira devolvida pela partição. & Exclui a posição final do pivô. \\
    Observação & Pode fazer menos trocas em certos dados. & É simples de explicar, mas pode trocar muitas vezes e degradar com chaves iguais. \\
    \bottomrule
  \end{tabularx}
  \vspace{.35em}
  \begin{beamercolorbox}[wd=\linewidth,sep=.45em]{block body}
    \scriptsize\textbf{Na \texttt{lesson11}:} o pivô é \texttt{values[right - 1]}; \texttt{store} avança para valores estritamente menores e, ao fim, recebe o pivô. É uma partição do tipo Lomuto.
  \end{beamercolorbox}
  \vfill
  \citep{cormen2022,sedgwick2011}
\end{frame}

\begin{frame}{Passo a passo (1/2): seis números, pivô 3}
  \small
  Entrada: \([5,\,2,\,9,\,1,\,6,\,3]\). Índices começam em zero; pivô é o último item.
  \vspace{.3em}
  \begin{tabularx}{\textwidth}{c c c X}
    \toprule
    \(j\) & Valor & Teste & Vetor após o passo \\
    \midrule
    -- & -- & Estado inicial; \texttt{store}=0 & \([5,2,9,1,6,3]\) \\
    0 & 5 & \(5<3\)? não & \([5,2,9,1,6,3]\) \\
    1 & 2 & \(2<3\)? sim; troca \(A[0]\leftrightarrow A[1]\) & \([2,5,9,1,6,3]\) \\
    2 & 9 & \(9<3\)? não & \([2,5,9,1,6,3]\) \\
    3 & 1 & \(1<3\)? sim; troca \(A[1]\leftrightarrow A[3]\) & \([2,1,9,5,6,3]\) \\
    4 & 6 & \(6<3\)? não & \([2,1,9,5,6,3]\) \\
    \bottomrule
  \end{tabularx}
  \vspace{.55em}
  \begin{center}
    \texttt{store}=2; valores menores que 3 estão à esquerda da fronteira. A varredura terminou.
  \end{center}
\end{frame}

\begin{frame}{Passo a passo (2/2): fixar o pivô e recursar}
  \small
  \begin{enumerate}
    \item Trocar \(A[\texttt{store}]\) com o pivô: \([2,1,3,5,6,9]\). O 3 está na posição final 2.
    \item Recursar à esquerda: \([2,1]\), pivô 1 \(\rightarrow [1,2]\).
    \item Recursar à direita: \([5,6,9]\), pivô 9 \(\rightarrow [5,6,9]\).
    \item Juntar não é necessário: a concatenação dos intervalos e do pivô já dá a ordem.
  \end{enumerate}
  \vspace{.5em}
  \begin{center}
    \begin{tikzpicture}[node distance=.35cm]
      \node[draw=MondrianBlack,fill=SoftGray,inner sep=6pt] (root) {\([5,2,9,1,6,3]\)};
      \node[draw=MondrianBlue,inner sep=5pt,below left=of root] (left) {\([2,1]\)};
      \node[draw=MondrianRed,inner sep=5pt,below right=of root] (right) {\([5,6,9]\)};
      \draw[-{Stealth},MondrianBlack] (root) -- (left);
      \draw[-{Stealth},MondrianBlack] (root) -- (right);
      \node[draw=MondrianBlack,fill=SoftGray,inner sep=5pt,below=1.15cm of root] (done) {\([1,2,3,5,6,9]\)};
      \draw[-{Stealth},MondrianBlack] (left) -- (done);
      \draw[-{Stealth},MondrianBlack] (right) -- (done);
    \end{tikzpicture}
  \end{center}
  \small Esta execução ilustra o rastreamento do código usado na implementação, não todas as variantes de Quick Sort.
\end{frame}

\begin{frame}{Pseudocódigo: partição Lomuto e Quick Sort}
  \small
  \begin{columns}[T,onlytextwidth]
    \column{.49\textwidth}
      \textbf{Particionar o intervalo \([lo,hi)\)}
      \begin{enumerate}
        \item \(p \gets A[hi-1]\); \(s \gets lo\)
        \item Para \(j=lo,\ldots,hi-2\):
        \begin{itemize}
          \item se \(A[j]<p\), trocar \(A[j]\) com \(A[s]\); \(s\gets s+1\)
        \end{itemize}
        \item Trocar \(A[s]\) com \(A[hi-1]\)
        \item Devolver \(s\)
      \end{enumerate}
    \column{.49\textwidth}
      \textbf{Ordenar}
      \begin{enumerate}
        \item Se \(hi-lo<2\), retornar.
        \item \(q\gets\operatorname{Particionar}(A,lo,hi)\)
        \item Ordenar \(A[lo..q)\)
        \item Ordenar \(A[q+1..hi)\)
      \end{enumerate}
      \vspace{.4em}
      \softbox{Invariante}{Durante a varredura, \(A[lo..s-1]<p\); os elementos ainda não examinados ficam à direita do índice \(j\).}
  \end{columns}
  \vfill
  \citep{cormen2022}
\end{frame}

\begin{frame}[fragile]{C: implementação da partição}
  \begin{lstlisting}[style=code]
static size_t partition(int *a, size_t lo, size_t hi)
{
    int pivot = a[hi - 1];
    size_t store = lo;

    for (size_t j = lo; j + 1 < hi; ++j) {
        if (a[j] < pivot) {
            int tmp = a[j];
            a[j] = a[store];
            a[store] = tmp;
            ++store;
        }
    }

    int tmp = a[store];
    a[store] = a[hi - 1];
    a[hi - 1] = tmp;
    return store;
}
  \end{lstlisting}
  \vspace{.05em}
  \scriptsize Pivô no último índice do intervalo semiaberto \([lo,hi)\); a partição devolve sua posição definitiva.
\end{frame}

\begin{frame}[fragile]{C: recursão sobre os intervalos}
  \begin{lstlisting}[style=code]
static void quick_sort(int *a, size_t lo, size_t hi)
{
    if (hi - lo < 2) return;
    size_t q = partition(a, lo, hi);
    quick_sort(a, lo, q);
    quick_sort(a, q + 1, hi);
}
  \end{lstlisting}
  \vspace{-.2em}
  \small A chamada inicial para \(n\) elementos é \texttt{quick\_sort(a, 0, n)}. O código da \texttt{lesson11} acrescenta métricas e atualização de progresso.
\end{frame}

\begin{frame}{Propriedades: o que o Quick Sort oferece (e não oferece)}
  \begin{columns}[T,onlytextwidth]
    \column{.48\textwidth}
      \softbox{Pontos fortes}{\begin{itemize}
        \item Bom desempenho médio: \(\Theta(n\log n)\).
        \item Particionamento in-place, com boa localidade de acesso.
        \item Laço interno compacto; baixo custo auxiliar fora da pilha.
        \item Pode ser adaptado a distintos tipos de dados e padrões.
      \end{itemize}}
    \column{.48\textwidth}
      \softbox{Limitações}{\begin{itemize}
        \item Pior caso \(\Theta(n^2)\) com partições repetidamente desequilibradas.
        \item Versão recursiva pode usar \(O(n)\) de pilha.
        \item Não é estável em sua forma usual.
        \item Resultado depende da política de pivô, da partição e de chaves repetidas.
      \end{itemize}}
  \end{columns}
  \vspace{.6em}
  \begin{center}
    \small     ``In-place'' não significa memória total \(O(1)\): a recursão também ocupa espaço.
  \end{center}
  \citep{cormen2022,knuth1998,sedgwick2011}
\end{frame}

\begin{frame}{Quando o pivô dá errado?}
  \begin{columns}[T,onlytextwidth]
    \column{.53\textwidth}
      Considere \(1,2,3,4,5\) com o maior elemento sempre escolhido como pivô:
      \[
      4+3+2+1=\frac{n(n-1)}2
      \]
      comparações no caso extremo.
      \begin{itemize}
        \item Pivô fixo pode ser vulnerável a entradas já ordenadas ou construídas para a implementação.
        \item Entradas com muitas chaves iguais exigem tratamento consciente.
      \end{itemize}
    \column{.43\textwidth}
      \begin{tikzpicture}[x=.61cm,y=.48cm]
        \foreach \y/\len in {3/5,2/4,1/3,0/2} {
          \draw[MondrianBlack] (0,\y) rectangle (\len,\y+.35);
          \fill[MondrianRed] (\len-.38,\y) rectangle (\len,\y+.35);
        }
        \node[font=\scriptsize,align=left] at (2.25,-.65) {Um lado quase vazio\\a cada partição};
      \end{tikzpicture}
  \end{columns}
  \vfill
  \citep{cormen2022,nist}
\end{frame}

\begin{frame}{Chaves repetidas: partição em três vias}
  \small
  Separar em três regiões: \(<p\), \(=p\) e \(>p\). A região igual não precisa ser recursada.
  \begin{center}
    \begin{tikzpicture}[x=.76cm,y=.75cm]
      \foreach \x/\v/\c in {0/2/MondrianBlue,1/2/MondrianBlue,2/4/MondrianYellow,3/4/MondrianYellow,4/4/MondrianYellow,5/7/MondrianRed,6/9/MondrianRed} {
        \draw[MondrianBlack] (\x,0) rectangle ++(.86,.68);
        \node at (\x+.43,.34) {\v};
      }
      \draw[decorate,decoration={brace,amplitude=4pt},MondrianBlack] (0,-.12) -- (1.86,-.12);
      \draw[decorate,decoration={brace,amplitude=4pt},MondrianBlack] (2,-.12) -- (4.86,-.12);
      \draw[decorate,decoration={brace,amplitude=4pt},MondrianBlack] (5,-.12) -- (6.86,-.12);
      \node[font=\scriptsize] at (.93,-.62) {\(<p\)};
      \node[font=\scriptsize] at (3.43,-.62) {\(=p\)};
      \node[font=\scriptsize] at (5.93,-.62) {\(>p\)};
    \end{tikzpicture}
  \end{center}
  \begin{itemize}
    \item A solução de Dijkstra para a ``bandeira holandesa'' mantém fronteiras para as três regiões.
    \item Em dados com muitas duplicatas, essa técnica evita recursão inútil sobre valores iguais; Bentley--McIlroy a incorporaram a uma implementação prática de ordenação.
    \item A implementação da \texttt{lesson11} não é tri-partida: usa comparação estrita \texttt{< pivot}.
  \end{itemize}
  \citep{dijkstra1976,bentley1993,sedgwick2011}
\end{frame}

\begin{frame}{Melhorias modernas: Quick Sort como componente}
  \small
  \begin{tabularx}{\textwidth}{>{\bfseries}p{3.1cm} X}
    \toprule
    Técnica & O que busca resolver \\
    \midrule
    Pivô aleatório & Reduz a dependência de uma entrada adversa fixa; expectativa \(O(n\log n)\). \\
    Mediana de três / amostragem & Evita pivôs extremos em padrões comuns, sem garantir o pior caso. \\
    Partição em três vias & Trata chaves repetidas com uma única região de igualdade. \\
    Inserção em intervalos pequenos & Diminui custos constantes e chamadas recursivas curtas. \\
    Introsort & Monitora a profundidade e recorre ao Heapsort para limitar o pior caso a \(O(n\log n)\). \\
    Dual-pivot / variantes adaptativas & Exploram políticas diferentes de partição e padrões dos dados; desempenho depende da implementação. \\
    \bottomrule
  \end{tabularx}
  \vspace{.55em}
  \small
  Exemplo de uso em biblioteca: a documentação de Java descreve Dual-Pivot Quicksort para certas ordenações de arrays primitivos. A estratégia concreta de \texttt{std::sort} depende da biblioteca; não se deve inferir uma implementação a partir apenas do nome da API.
  \citep{musser1997,bentley1993,javaarrays}
\end{frame}

\begin{frame}{Comparação com outros algoritmos do trabalho}
  \scriptsize
  \begin{tabularx}{\textwidth}{>{\bfseries}p{2.05cm} p{2.05cm} p{1.25cm} p{1.35cm} X}
    \toprule
    Algoritmo & Tempo típico / pior & Estável? & In-place? & Observação principal \\
    \midrule
    Quick Sort & \(n\log n\) / \(n^2\) & Não & Em geral & Rápido na prática; o pivô importa. \\
    Merge Sort & \(n\log n\) / \(n\log n\) & Sim* & Não* & Garantia de pior caso; memória auxiliar. \\
    Heap Sort & \(n\log n\) / \(n\log n\) & Não & Sim & Pior caso garantido; acesso menos sequencial. \\
    Insertion Sort & \(n^2\) / \(n^2\) & Sim & Sim & Simples e eficiente para entradas pequenas/quase ordenadas. \\
    Selection Sort & \(n^2\) / \(n^2\) & Não* & Sim & Poucas trocas; muitas comparações. \\
    Bubble Sort & \(n^2\) / \(n^2\) & Sim* & Sim & Didático; depende da otimização para parar cedo. \\
    Radix Sort & \(d(n+b)\) / \(d(n+b)\) & Depende & Não* & Não compara chaves; custo depende da representação e base. \\
    \bottomrule
  \end{tabularx}
  \vspace{.4em}
  \tiny \(*\) Propriedades da versão usual/implementação estável, não de todo algoritmo que compartilha o nome. \(d\): número de dígitos/passagens; \(b\): base ou tamanho do alfabeto.
  \vfill
  \citep{cormen2022,ziviani2011,knuth1998}
\end{frame}

\begin{frame}{Busca: o que muda quando os dados estão ordenados?}
  \begin{columns}[T,onlytextwidth]
    \column{.48\textwidth}
      \softbox{Busca sequencial}{Examina itens até encontrar a chave ou chegar ao fim.\\[.3em]
      Pior caso: \(O(n)\). Não exige ordenação.}
    \column{.48\textwidth}
      \softbox{Busca binária}{Compara com o elemento central e descarta metade do intervalo a cada etapa.\\[.3em]
      Pior caso: \(O(\log n)\). Exige sequência ordenada e acesso indexado eficiente.}
  \end{columns}
  \vspace{.65em}
  \begin{center}
    \small
    Na \texttt{lesson11}, o benchmark de buscas usa 10.000 chaves e 10.000 consultas pseudoaleatórias: sequencial no vetor original e binária em uma cópia ordenada.
  \end{center}
  \citep{cormen2022}
\end{frame}

\begin{frame}{Metodologia experimental da \texttt{lesson11}}
  \begin{itemize}
    \item O módulo executa cada algoritmo sobre arquivos de entrada de \(100\), \(1.000\), \(10.000\) e \(100.000\) inteiros.
    \item Mede o intervalo da chamada ao algoritmo com contador de alta resolução do Windows; o tempo de gravação do arquivo é posterior à medição.
    \item Os arquivos ordenados preservam os valores da entrada. A checagem local dos 21 resultados disponíveis confirmou ordem não decrescente e mesmo multiconjunto.
    \item Os nomes dos arquivos guardam o tempo arredondado para microssegundos inteiros ou milissegundos inteiros; os arquivos não registram máquina, compilador, repetições ou dispersão.
  \end{itemize}
  \vspace{.45em}
  \softbox{Escopo observado}{Há artefatos de dados para sete ordenações até \(10.000\) elementos. Para \(100.000\), há resultados de Heap, Merge, Quick e Radix; não há arquivos dessa dimensão para Bubble, Insertion e Selection.}
  \vfill
  \tiny Fonte: \texttt{problems/lesson11/problem\_110.c}, \texttt{search-benchmark.c} e \texttt{data/}.
\end{frame}

\begin{frame}{Expectativa teórica: crescimento, não previsão de milissegundos}
  \centering
  \begin{tikzpicture}
    \begin{axis}[
      width=.84\textwidth,height=5.1cm,
      xmode=log,ymode=log,
      log basis x=10,log basis y=10,
      xlabel={Tamanho \(n\) (escala logarítmica)},
      ylabel={Operações normalizadas (escala logarítmica)},
      xmin=100,xmax=100000,ymin=.7,ymax=2000000,
      xtick={100,1000,10000,100000},
      xticklabels={100,1.000,10.000,100.000},
      ymajorgrids=true,xmajorgrids=true,
      grid style={draw=SoftGray},
      legend style={at={(.02,.98)},anchor=north west,draw=none,fill=white,font=\scriptsize},
      tick label style={font=\scriptsize},
      label style={font=\scriptsize}
    ]
      \addplot[very thick,MondrianBlue,mark=square*] coordinates {(100,1) (1000,15.05) (10000,200.1) (100000,250.1)};
      \addlegendentry{\(n\log_2 n\), normalizado em \(n=100\)}
      \addplot[very thick,MondrianRed,mark=*] coordinates {(100,1) (1000,100) (10000,10000) (100000,1000000)};
      \addlegendentry{\(n^2\), normalizado em \(n=100\)}
    \end{axis}
  \end{tikzpicture}
  \vspace{-.1em}
  \small Curvas assintóticas ilustrativas: Quick Sort balanceado segue a tendência \(n\log n\); partições extremas seguem \(n^2\). Não são tempos medidos.
\end{frame}

\begin{frame}{Resultados práticos registrados: crescimento por tamanho}
  \centering
  \begin{tikzpicture}
    \begin{axis}[
      width=.88\textwidth,height=5.5cm,
      ymode=log,log basis y=10,
      xlabel={Número de elementos},
      ylabel={Tempo codificado no nome do arquivo (ms; log)},
      xmin=100,xmax=100000,ymin=.004,ymax=200,
      xtick={100,1000,10000,100000},
      xticklabels={100,1.000,10.000,100.000},
      ymajorgrids=true,xmajorgrids=true,
      grid style={draw=SoftGray},
      legend columns=4,
      legend style={at={(.5,-.22)},anchor=north,draw=none,font=\tiny},
      tick label style={font=\scriptsize},
      label style={font=\scriptsize},
      clip=false
    ]
      \addplot[mark=square*,MondrianRed,thick] coordinates {(100,.022) (1000,2) (10000,117)};
      \addlegendentry{Bubble}
      \addplot[mark=triangle*,MondrianBlue,thick] coordinates {(100,.006) (1000,.091) (10000,1) (100000,13)};
      \addlegendentry{Heap}
      \addplot[mark=*,MondrianBlack,thick] coordinates {(100,.008) (1000,.709) (10000,48)};
      \addlegendentry{Insertion}
      \addplot[mark=diamond*,MondrianYellow,thick] coordinates {(100,.008) (1000,.090) (10000,.920) (100000,11)};
      \addlegendentry{Merge}
      \addplot[mark=pentagon*,MondrianBlue,dashed,thick] coordinates {(100,.005) (1000,.064) (10000,.750) (100000,10)};
      \addlegendentry{Quick}
      \addplot[mark=oplus,MondrianRed,dashed,thick] coordinates {(100,.012) (1000,.110) (10000,1) (100000,9)};
      \addlegendentry{Radix}
      \addplot[mark=x,MondrianBlack,dashed,thick] coordinates {(100,.013) (1000,1) (10000,90)};
      \addlegendentry{Selection}
    \end{axis}
  \end{tikzpicture}
  \vspace{-.1em}
  \tiny Tempos aproximados a partir dos sufixos dos arquivos \texttt{data/*-sort/*}. Os pontos de \(100.000\) estão disponíveis somente para Heap, Merge, Quick e Radix; o gráfico não interpola dados ausentes.
\end{frame}

\begin{frame}{Curvas teóricas versus medidas: escala relativa}
  \scriptsize
  \begin{columns}[T,onlytextwidth]
    \column{.48\textwidth}
      \textbf{Previsão assintótica normalizada}
      \begin{tabularx}{\linewidth}{>{\bfseries}l r r r r}
        \toprule
        Modelo & 100 & 1.000 & 10.000 & 100.000 \\
        \midrule
        \(n\log_2 n\) & 1 & 15,1 & 200,1 & 2.501 \\
        \(n^2\) & 1 & 100 & 10.000 & 1.000.000 \\
        \(n\) & 1 & 10 & 100 & 1.000 \\
        \bottomrule
      \end{tabularx}
      \vspace{.35em}
      \tiny A normalização elimina a constante desconhecida e permite comparar a forma do crescimento.
    \column{.48\textwidth}
      \textbf{Tempos medidos normalizados}
      \begin{tabularx}{\linewidth}{>{\bfseries}l r r r r}
        \toprule
        Algoritmo & 100 & 1.000 & 10.000 & 100.000 \\
        \midrule
        Quick & 1 & 12,8 & 150 & 2.000 \\
        Heap & 1 & 15,2 & 166,7 & 2.166,7 \\
        Merge & 1 & 11,3 & 115 & 1.375 \\
        Radix & 1 & 9,2 & 83,3 & 750 \\
        \bottomrule
      \end{tabularx}
      \vspace{.35em}
      \tiny Base: o menor arquivo de cada algoritmo; valores aproximados pelos nomes dos arquivos.
  \end{columns}
  \vspace{.5em}
  \begin{beamercolorbox}[wd=\linewidth,sep=.55em]{block body}
    \textbf{Leitura:} Quick, Heap e Merge crescem na mesma ordem de grandeza esperada para \(n\log n\), embora suas constantes e arredondamentos sejam diferentes. Radix apresenta crescimento próximo de linear nesta faixa, coerente com \(d\) e a base fixados, mas não é diretamente comparável a um algoritmo baseado em comparações.
  \end{beamercolorbox}
\end{frame}

\begin{frame}{Comparação por algoritmo: previsão e evidência}
  \scriptsize
  \begin{tabularx}{\textwidth}{>{\bfseries}p{1.7cm} p{2.1cm} p{2.5cm} X}
    \toprule
    Algoritmo & Previsão teórica & Curva observada & Conclusão para esta execução \\
    \midrule
    Quick Sort & Médio \(n\log n\); pior \(n^2\) & \(0,005 \to 10\) ms; fator \(\approx 2.000\) & A série observada acompanha a família \(n\log n\); pivô final e entrada aleatória não produziram degradação quadrática. \\
    Heap Sort & \(n\log n\) em todos os casos & \(0,006 \to 13\) ms; fator \(\approx 2.167\) & Crescimento próximo ao de Quick/Merge; oferece garantia de pior caso. \\
    Merge Sort & \(n\log n\) em todos os casos & \(0,008 \to 11\) ms; fator \(\approx 1.375\) & Compatível com \(n\log n\), com constante favorável nesta medição. \\
    Radix Sort & \(d(n+b)\); linear se \(d,b\) fixos & \(0,012 \to 9\) ms; fator \(\approx 750\) & Crescimento mais próximo de \(n\) na faixa; depende do domínio numérico e da implementação. \\
    Insertion Sort & \(n^2\) no pior/médio caso & \(0,008 \to 48\) ms em 10 mil; 100 mil ausente & A forte subida entre 1 mil e 10 mil é compatível com custo quadrático; não extrapolar como medição. \\
    Selection Sort & \(n^2\) & \(0,013 \to 79\) ms em 10 mil; 100 mil ausente & Tendência quadrática aparente; a ausência de 100 mil impede afirmar o tempo real. \\
    Bubble Sort & \(n^2\) (versão usual) & \(0,022 \to 117\) ms em 10 mil; 100 mil ausente & Maior custo entre os quadráticos medidos; não há curva observada em 100 mil. \\
    \bottomrule
  \end{tabularx}
  \vspace{.1em}
  \tiny ``Compatível'' significa que os pontos não contradizem a ordem de crescimento esperada; não é uma prova estatística da complexidade.
\end{frame}

\begin{frame}{Como interpretar as curvas sem confundir teoria e prática}
  \small
  \begin{enumerate}
    \item \textbf{Teoria:} compara ordens de crescimento quando \(n\) aumenta; ignora constantes, hardware, cache e implementação.
    \item \textbf{Prática:} os nomes dos arquivos registram uma execução arredondada; pontos próximos podem parecer iguais por causa da resolução.
    \item \textbf{Curvas paralelas:} Quick, Heap e Merge mantêm crescimento semelhante no gráfico de 100 a 100.000.
    \item \textbf{Curvas ausentes:} não há ponto de 100.000 para Bubble, Insertion e Selection; uma linha não deve ser inventada por extrapolação.
    \item \textbf{Radix:} sua previsão depende de \(d\) e \(b\); a comparação com \(n\log n\) é uma comparação de modelos, não de operações equivalentes.
  \end{enumerate}
  \vspace{.15em}
  \softbox{Mensagem para a banca}{Os dados são coerentes com as previsões assintóticas nas entradas observadas, mas a conclusão deve ser limitada ao ambiente, às entradas e aos arquivos efetivamente medidos.}
\end{frame}

\begin{frame}{Leitura dos resultados: o que os dados sustentam}
  \scriptsize
  \begin{tabularx}{\textwidth}{>{\bfseries}p{2.05cm} c c c c X}
    \toprule
    Algoritmo & \(n=100\) & \(n=1.000\) & \(n=10.000\) & \(n=100.000\) & Leitura cuidadosa \\
    \midrule
    Quick Sort & 5 \(\mu s\) & 64 \(\mu s\) & 750 \(\mu s\) & 10 ms & Menor tempo registrado nos quatro tamanhos disponíveis. \\
    Heap Sort & 6 \(\mu s\) & 91 \(\mu s\) & 1 ms & 13 ms & Crescimento compatível com \(n\log n\). \\
    Merge Sort & 8 \(\mu s\) & 90 \(\mu s\) & 920 \(\mu s\) & 11 ms & Mesmo patamar assintótico de Quick e Heap. \\
    Insertion Sort & 8 \(\mu s\) & 709 \(\mu s\) & 48 ms & -- & Não há arquivo de 100 mil. \\
    Selection Sort & 13 \(\mu s\) & 1 ms & 79 ms & -- & Não há arquivo de 100 mil; custo quadrático sugere crescimento elevado. \\
    Bubble Sort & 22 \(\mu s\) & 1 ms & 117 ms & -- & Não há arquivo de 100 mil; custo quadrático sugere crescimento elevado. \\
    Radix Sort & 12 \(\mu s\) & 110 \(\mu s\) & 1 ms & 9 ms & Melhor resultado registrado em 100 mil nesta amostra. \\
    \bottomrule
  \end{tabularx}
  \vspace{.1em}
  \tiny Nomes consultados: \texttt{100000\_13ms.txt}, \texttt{100000\_11ms.txt}, \texttt{100000\_10ms.txt} e \texttt{100000\_9ms.txt}; tempos arredondados pela rotina de gravação.
\end{frame}

\begin{frame}{O que não se pode concluir apenas com esses tempos}
  \begin{itemize}
    \item Uma execução por algoritmo/tamanho não estima variabilidade nem intervalo de confiança.
    \item Resolução diferente entre microssegundos e milissegundos esconde diferenças em alguns pontos.
    \item Não foram registrados aqui CPU, sistema, compilador, flags, estado de cache ou distribuição detalhada das entradas.
    \item Quick Sort é o menor tempo \textit{registrado} nesta amostra, mas isso não prova superioridade universal.
  \end{itemize}
  \vspace{.4em}
  \softbox{Como fortalecer o experimento}{Executar várias repetições; usar as mesmas entradas e cópias independentes; relatar mediana e dispersão; registrar compilador/flags e hardware; testar aleatória, ordenada, reversa e muitos valores iguais; completar os três algoritmos quadráticos em \(100.000\) somente se o custo for aceitável.}
\end{frame}

\begin{frame}{Por que ainda escolher Quick Sort?}
  \begin{columns}[T,onlytextwidth]
    \column{.48\textwidth}
      \softbox{É uma boa escolha quando}{\begin{itemize}
        \item dados cabem em memória;
        \item estabilidade não é requisito;
        \item desempenho médio e localidade são prioritários;
        \item a implementação pode escolher boas políticas de pivô e fallback.
      \end{itemize}}
    \column{.48\textwidth}
      \softbox{Considere alternativas quando}{\begin{itemize}
        \item pior caso garantido é obrigatório (Heap/Merge/Introsort);
        \item estabilidade é obrigatória (Merge);
        \item dados são quase ordenados e pequenos (Insertion);
        \item a chave tem estrutura explorável (Radix).
      \end{itemize}}
  \end{columns}
  \vspace{.5em}
  \begin{center}
    \small A escolha é uma decisão de engenharia: modelo de dados, requisitos e implementação importam tanto quanto a notação assintótica.
  \end{center}
  \citep{knuth1998,cormen2022,ziviani2011}
\end{frame}

\begin{frame}{Conexão com o trabalho da disciplina}
  \begin{columns}[T,onlytextwidth]
    \column{.57\textwidth}
      \begin{itemize}
        \item A implementação integra o menu de comparação da \texttt{lesson11}.
        \item A rotina acompanha comparações, movimentos e progresso.
        \item Arquivos ordenados são gravados em \texttt{data/quick-sort/}.
        \item Os artefatos disponíveis cobrem \(100\), \(1.000\), \(10.000\) e, parcialmente, \(100.000\) números; foram validados quanto à ordenação e preservação dos valores.
      \end{itemize}
      \vspace{.35em}
      \scriptsize Fontes: \texttt{problems/lesson11/quick-sort.c}, \texttt{problem\_110.c} e \texttt{data/quick-sort/}.
    \column{.4\textwidth}
      \centering
      \placeholder{5.2cm}{3.0cm}\\[-.2em]
      \scriptsize Substituir por uma captura da tela de comparação ou do menu da \texttt{lesson11}.
  \end{columns}
\end{frame}

\begin{frame}{Síntese}
  \begin{enumerate}
    \item Quick Sort ordena por partições sucessivas e é um exemplo clássico de dividir e conquistar.
    \item A qualidade do particionamento determina a forma da recursão: \(n\log n\) quando equilibrada; \(n^2\) no pior caso.
    \item A implementação da \texttt{lesson11} usa pivô final e partição Lomuto-style; duplicatas e pivôs extremos são pontos de atenção.
    \item Variantes modernas combinam partições melhores com inserção, três vias ou fallback para limitar riscos.
    \item Os resultados disponíveis favorecem Quick Sort nesta amostra, mas a conclusão deve respeitar resolução e limites experimentais.
  \end{enumerate}
  \vspace{.5em}
  \begin{center}
    \Large\bfseries Obrigado(a)!\\[-.1em]
    \normalsize Perguntas?
  \end{center}
\end{frame}

\begin{frame}[allowframebreaks]{Referências}
  \footnotesize
  \begin{thebibliography}{99}
    \bibitem{hoare1962}
      HOARE, C. A. R. Quicksort. \textit{The Computer Journal}, v. 5, n. 1, p. 10--16, 1962.
      \href{https://doi.org/10.1093/comjnl/5.1.10}{doi:10.1093/comjnl/5.1.10}.

    \bibitem{cormen2022}
      CORMEN, T. H.; LEISERSON, C. E.; RIVEST, R. L.; STEIN, C.
      \textit{Introduction to Algorithms}. 4. ed. MIT Press, 2022.

    \bibitem{knuth1998}
      KNUTH, D. E. \textit{The Art of Computer Programming, Volume 3: Sorting and Searching}.
      2. ed. Addison-Wesley, 1998.

    \bibitem{sedgwick2011}
      SEDGEWICK, R.; WAYNE, K. \textit{Algorithms}. 4. ed. Addison-Wesley, 2011.
      Material de apoio: \url{https://algs4.cs.princeton.edu/23quicksort/}.

    \bibitem{dijkstra1976}
      DIJKSTRA, E. W. \textit{A Discipline of Programming}. Prentice-Hall, 1976.

    \bibitem{bentley1993}
      BENTLEY, J. L.; McILROY, M. D. Engineering a sort function.
      \textit{Software: Practice and Experience}, v. 23, n. 11, p. 1249--1265, 1993.
      \href{https://doi.org/10.1002/spe.4380231105}{doi:10.1002/spe.4380231105}.

    \bibitem{musser1997}
      MUSSER, D. R. Introspective sorting and selection algorithms.
      \textit{Software: Practice and Experience}, v. 27, n. 8, p. 983--993, 1997.
      \href{https://doi.org/10.1002/(SICI)1097-024X(199708)27:8\%3C983::AID-SPE117\%3E3.0.CO;2-\%23}
      {doi:10.1002/(SICI)1097-024X(199708)27:8<983::AID-SPE117>3.0.CO;2-\#}.

    \bibitem{ziviani2011}
      ZIVIANI, N. \textit{Projeto de Algoritmos: com implementações em Pascal e C}.
      3. ed. Cengage Learning, 2011.

    \bibitem{nist}
      MARTINEZ, C. Quicksort. In: BLACK, P. E. (ed.). \textit{Dictionary of Algorithms and Data Structures}.
      NIST, 2023. Disponível em: \url{https://xlinux.nist.gov/dads/HTML/quicksort.html}.

    \bibitem{javaarrays}
      ORACLE. \textit{Java Platform, Standard Edition 8 API Specification: Arrays}.
      Nota de implementação sobre Dual-Pivot Quicksort.
      \url{https://docs.oracle.com/javase/8/docs/api/java/util/Arrays.html}.

    \bibitem{lesson11}
      CARAUMA, J. A. Repositório \textit{art-mondrian}, módulo \texttt{problems/lesson11}.
      Código e dados experimentais da implementação. Disponível em:
      \url{https://github.com/carauma7/art-mondrian}.
  \end{thebibliography}
\end{frame}

\end{document}
